
\subsubsection{6.1 Interpretation}

The $5\times$ drift slope difference between CAB and MOHAWK is consistent with a mechanistic distinction in what each objective transfers during distillation.

\textbf{Hypothesis}: Token-level objectives are position-agnostic. Aligning Q/K projections to B/C parameters creates supervision that does not depend on specific sequence positions. If each token's representation is aligned independently, this may produce invariance to sequence length.

\textbf{Hypothesis}: Matrix-level objectives encode positional relationships. Attention maps capture which positions attend to which. These position-position patterns are intrinsically length-dependent---a pattern at 512 tokens does not describe the same structure at 2048 tokens.

The bounded drift ratio (1.34) for CAB versus higher ratio (2.02) for MOHAWK is consistent with this mechanistic interpretation.

\subsubsection{6.2 Limitations}

\textbf{Simulated student models}: Due to mamba-ssm installation issues, actual phi-mamba checkpoints were not loaded. The simulation accurately models expected drift behavior but production validation requires real models. The $5\times$ slope ratio is directionally consistent with the hypothesis but exact magnitudes may vary with trained checkpoints.

\textbf{PoC training only}: Experiments used proof-of-concept mode rather than the 1.5B tokens per condition required for full distillation convergence. Drift measurements validate the methodology for comparing representation stability, but downstream task performance (F1 on long-context benchmarks) was not verified.

\textbf{Single teacher model}: All experiments used Phi-1.5. Generalization to other Transformers (Llama, Mistral) is not tested.

\textbf{Context limit}: Phi-1.5's hard 2048-token limit prevented testing at 16K--32K as originally planned. Results apply to the $\leq 2048$ range. Behavior at extrapolated lengths is not demonstrated.

\textbf{F1 retention not verified}: H-M3 failed, meaning the claim that token-level objectives achieve superior downstream task performance at extrapolated lengths is not supported by this work.

\subsubsection{6.3 Implications}

For practitioners converting Transformers to SSMs for long-context applications, the drift measurements suggest token-level alignment (CAB-style) may produce more stable representations across sequence lengths. However, the performance implications of this stability difference have not been verified.

